\section{Síntesis y ampliación}

\subsection{Panorama de las ideas centrales}
\begin{table}[H]
\centering
\begin{tabular}{p{2.8cm}p{4.0cm}p{5.0cm}}
\toprule
\textbf{Tema} & \textbf{Conclusión del panel} & \textbf{Problema de sistemas correspondiente} \\
\midrule
Infraestructura de inferencia & La separación PD y el EP grande ya son el estándar en producción & ¿Cómo reorganizar los recursos según la carga de trabajo? \\
Infraestructura de RL & El rollout es el cuello de botella central; combinar entrenamiento e inferencia es cada vez más importante & ¿Cómo hacer que el RL de modelos grandes sea escalable? \\
Infraestructura de agentes & El entorno, el workflow y los flujos de eventos hacen el sistema más dinámico & ¿Cómo dar cabida a herramientas reales y fallos reales? \\
Sistema de entrenamiento & Se pasa de un framework de una sola línea principal a un sistema de programación heterogéneo & ¿Cómo separar en módulos los problemas de infraestructura? \\
Dirección futura & No se trata solo de entrenar el modelo, sino de entrenar el sistema & ¿Cómo hacer que el sistema y el algoritmo iteren juntos? \\
\bottomrule
\end{tabular}
\caption{Las cinco líneas más importantes que recordar del panel de infraestructura}
\end{table}

\subsection{Tres conclusiones orientadas a la práctica}
\begin{enumerate}
  \item Si la carga de trabajo ya se volvió Agent / RL / multimodal, no se debe seguir diseñando el sistema con la mirada de un preentrenamiento puro.
  \item Lo que cada vez más limita el techo de las capacidades de los modelos grandes es el rollout, el entorno, la programación y la recuperación de fallos, no la velocidad de un solo operador.
  \item La señal de que una infraestructura es madura no es que el usuario perciba su complejidad, sino que se mantenga estable, reproducible y escalable bajo cargas de trabajo complejas.
\end{enumerate}

\begin{importantbox}{La verdadera enseñanza del panel de infraestructura}
Nos recuerda que, en la era de los modelos grandes, el ``sistema'' no es un accesorio, sino parte de la capacidad. Muchos problemas que parecen cuellos de botella algorítmicos solo pueden resolverse de verdad cuando el diseño del sistema, la construcción del entorno y las interfaces de despliegue avanzan al mismo ritmo.
\end{importantbox}

\subsection{Lecturas adicionales}
\begin{itemize}
  \item Material sobre sistemas de inferencia como MoonCake, KTransformers, vLLM y SGLang
  \item Proyectos de código abierto y artículos relacionados con la infraestructura de RL y los rollout engines
  \item Trabajos sobre entornos de agentes, modelos de mundo y construcción de benchmarks
  \item Experiencia de migración de entrenamiento distribuido, recuperación ante fallos, sistemas de archivos y bases de datos hacia la infraestructura de IA
  \item Prácticas de ingeniería sobre heterogeneous workflow y sistemas de IA basados en eventos
\end{itemize}

\end{document}
